\documentclass[12pt]{article}
\usepackage[a4paper, margin=1in]{geometry}
\usepackage{amsmath, amssymb}
\usepackage{graphicx}

\title{CSE400 -- Lecture 20 Scribe}
\author{}
\date{}

\begin{document}

\maketitle

\section*{Linear Transformations and Expectations of Multiple Random Variables}

\section*{1. Outline}
\begin{itemize}
\item Statistical Representation of Random Vectors
\item Mean Vector, Correlation, and Covariance Matrices
\item Linear Transformations
\item Mapping properties under Multivariate Gaussian Distributions
\item From 1D/2D cases to N-dimensional Joint PDFs
\item Special Cases and Properties
\item Implications of zero correlation and matrix factorization
\end{itemize}

\section*{2. Random Vector Representation}

\[
Y = AX + b
\]

\textbf{Definition}

For random variables $X_1, X_2, \dots, X_N$, we represent them as:

\[
X =
\begin{bmatrix}
X_1 \\
X_2 \\
\vdots \\
X_N
\end{bmatrix}
\]

$X$ is called a random vector.

Provides a compact representation of multiple random variables.

\section*{3. Mean Vector}

\textbf{Definition}

\[
\mu_X = E[X]
\]

\[
\mu_X =
\begin{bmatrix}
E[X_1] \\
E[X_2] \\
\vdots \\
E[X_N]
\end{bmatrix}
\]

Represents the average value of the random vector.

\section*{4. Correlation Matrix}

\textbf{Definition}

\[
R_{XX} = E[XX^T]
\]

Element-wise:
\[
R_{XX}(i,j) = E[X_i X_j]
\]

Matrix form:
\[
R_{XX} =
\begin{bmatrix}
E[X_1 X_1] & E[X_1 X_2] & \cdots \\
E[X_2 X_1] & E[X_2 X_2] & \cdots \\
\vdots & \vdots & \ddots
\end{bmatrix}
\]

$R_{XX}$ is symmetric.

\section*{5. Covariance Matrix}

\textbf{Definition}

\[
C_{XX} = E[(X - \mu_X)(X - \mu_X)^T]
\]

Structure:
\[
C_{XX} =
\begin{bmatrix}
\text{Var}(X_1) & \text{Cov}(X_2,X_1) \\
\text{Cov}(X_1,X_2) & \text{Var}(X_2)
\end{bmatrix}
\]

Symmetry:
\[
\text{Cov}(X_i,X_j) = \text{Cov}(X_j,X_i)
\]

\section*{6. Linear Transformation of Random Vectors}

\[
Y = AX + b
\]

$A$: transformation matrix (scaling, rotation) \\
$b$: constant vector (shift/translation)

Component form:
\[
Y_i = \sum_{j=1}^{N} a_{ij} X_j + b_i
\]

\section*{7. Mean of Transformed Vector}

\[
\mu_Y = E[Y]
\]

Step-by-step:

\[
\mu_Y = E[AX + b]
\]
\[
\mu_Y = E[AX] + E[b]
\]
\[
\mu_Y = A\mu_X + b
\]

\section*{8. Correlation Matrix Transformation}

\[
R_{YY} = E[YY^T] = E[(AX + b)(AX + b)^T]
\]

\[
R_{YY} = AE[XX^T]A^T + AE[X]b^T + bE[X]^T A^T + bb^T
\]

\[
R_{YY} = AR_{XX}A^T + A\mu_X b^T + b\mu_X^T A^T + bb^T
\]

\section*{9. Covariance Matrix Transformation}

\[
Y - \mu_Y = (AX + b) - (A\mu_X + b) = A(X - \mu_X)
\]

\[
C_{YY} = E[(Y - \mu_Y)(Y - \mu_Y)^T]
\]

\[
C_{YY} = E[A(X - \mu_X)(X - \mu_X)^T A^T]
\]

\[
C_{YY} = AC_{XX}A^T
\]

Important: Shift $b$ does not affect covariance.

\section*{10. Summary of Linear Transformation}

\begin{itemize}
\item Mean: $\mu_Y = A\mu_X + b$
\item Covariance: $C_{YY} = AC_{XX}A^T$
\item Correlation:
\[
R_{YY} = AR_{XX}A^T + A\mu_X b^T + b\mu_X^T A^T + bb^T
\]
\end{itemize}

\section*{11. Numerical Example}

Given $Y = AX$

(a) Mean:
\[
\mu_Y = A\mu_X
\]

(b) Correlation:
\[
R_{YY} = AR_{XX}A^T
\]

(c) Covariance:
\[
C_{YY} = R_{YY} - \mu_Y \mu_Y^T
\]

\section*{12. Gaussian Distributions}

1D Gaussian: Defined for a single random variable.

2D Joint Gaussian: Joint PDF given in slides.

\section*{13. Multivariate Gaussian Distribution}

For $N$-dimensional vector:

\[
X = [X_1, X_2, \dots, X_N]^T
\]

\[
f_X(x) = \frac{1}{(2\pi)^{N/2} \sqrt{\det(C_{XX})}} 
\exp\left(-\frac{1}{2}(x - \mu_X)^T C_{XX}^{-1} (x - \mu_X)\right)
\]

Components:
\begin{itemize}
\item $\mu_X$: mean vector
\item $C_{XX}$: covariance matrix
\end{itemize}

\section*{14. Special Case: Uncorrelated Gaussian Variables}

Condition:
\[
\text{Cov}(X_i, X_j) = 0 \quad \forall i \neq j
\]

Implication:
\begin{itemize}
\item Covariance matrix becomes diagonal
\item Joint PDF factorizes into product of individual Gaussians
\end{itemize}

\section*{15. Homework Problems}

1. Mean from matrices:
\[
\mu_X \mu_X^T = R_{XX} - C_{XX}
\]

2. Variance identity:
\[
\text{Var}[X + Y] = \text{Var}[X] + \text{Var}[Y] + 2\text{Cov}(X,Y)
\]

\section*{16. Case Study: Smart City Traffic + Weather System}

Step 1: Define random vector:
\[
X = [X_1, X_2, X_3, X_4, X_5]
\]

Where:
\begin{itemize}
\item $X_1$: Traffic flow
\item $X_2$: Speed
\item $X_3$: Rainfall
\item $X_4$: Temperature
\item $X_5$: AQI
\end{itemize}

Step 2: Mean vector:
\[
\mu_X = E[X]
\]

Step 3: Covariance matrix represents dependencies.

Examples:
\[
\text{Cov}(X_1, X_3) > 0
\]
\[
\text{Cov}(X_2, X_3) < 0
\]

Step 4: Gaussian assumption:
\[
X \sim N(\mu, \Sigma)
\]

Step 6: Linear transformation:
\[
Y = 0.6X_1 - 0.4X_2 + 0.8X_3
\]

\[
E[Y] = a\mu_X
\]

\section*{17. Key Insights}

\begin{itemize}
\item Random vector = system state
\item Covariance = relationships
\item Gaussian = structured uncertainty
\item Transformation = derived metrics
\end{itemize}

\section*{End of Lecture}

\end{document}